% =====================================================================
% NOTA (remover antes da entrega final):
% Este documento segue a estrutura exigida em formatorelatorios.pdf
% (artigo IEEE de conferencia, duas colunas, sem folha de capa, secoes a-e).
%
% PONTO A CONFIRMAR COM A PROFESSORA:
% O formato geral diz "max 03 autores", mas o grupo tem 4 integrantes.
% Para o tamanho, adotou-se o limite especifico de 4 paginas do roteiro.
% As secoes sobre topologia de rede foram adaptadas ao ambiente de execucao.
%
% Para compilar: pdflatex relatorio.tex (duas vezes, para o indice/refs)
% Se nao houver LaTeX instalado localmente, use o Overleaf
% (overleaf.com) e importe este arquivo com o template "IEEE Conference".
% =====================================================================
\documentclass[conference]{IEEEtran}
\usepackage[utf8]{inputenc}
\usepackage[T1]{fontenc}
\usepackage[brazilian]{babel}
\usepackage{amsmath}
\usepackage{graphicx}
\usepackage{listings}
\usepackage{xcolor}
\usepackage[hidelinks]{hyperref}

\newcommand{\pendente}[1]{\textcolor{red}{\textbf{[PENDENTE: #1]}}}

\begin{document}

\title{Cifra de Vigenère: Implementação e Criptoanálise\\
\large Trabalho de Implementação 1 --- CIC0201 Segurança Computacional}

% NOTA: formatorelatorios.pdf pede no maximo 3 autores; o grupo tem 4.
% Ajustar conforme orientacao da professora.
\author{
\IEEEauthorblockN{Ana Luísa Reis Nascente}
\IEEEauthorblockA{Matrícula: 211045688}
\and
\IEEEauthorblockN{Gabriel de Sousa}
\IEEEauthorblockA{Matrícula: 211056000}
\and
\IEEEauthorblockN{Marina Pimentel Moreno}
\IEEEauthorblockA{Matrícula: 222014071}
\and
\IEEEauthorblockN{Pedro de Paula Campos}
\IEEEauthorblockA{Matrícula: 231036050}
}

\maketitle

\begin{abstract}
Este trabalho apresenta a implementação da cifra de Vigenère e de um
ataque para recuperação da chave por propriedades estatísticas da
linguagem. A implementação atual inclui o cifrador e a estimativa do
tamanho da chave por Índice de Coincidência, validados por testes. A
recuperação das letras da chave e os resultados sobre os dois
criptogramas serão incorporados após a integração final.
\end{abstract}

\begin{IEEEkeywords}
cifra de Vigenère, criptografia clássica, criptoanálise, índice de coincidência, análise de frequência
\end{IEEEkeywords}

\section{Introdução}

Este relatório se insere no módulo de Criptografia Simétrica da
disciplina CIC0201 -- Segurança Computacional, e documenta o Trabalho
de Implementação 1: a cifra de Vigenère, cifra de substituição
polialfabética historicamente relevante por ilustrar, de forma
simples, dois princípios centrais da criptografia moderna --- a
diferença entre um espaço de chaves grande e segurança real, e como a
reutilização cíclica de uma chave curta introduz uma estrutura
estatística explorável \cite{roteiro}.

O trabalho está dividido em duas partes complementares: (i) um
cifrador/decifrador a ser implementado do zero, sem bibliotecas
criptográficas prontas; e (ii) um ataque de criptoanálise capaz de
recuperar a chave de um criptograma desconhecido a partir de
propriedades estatísticas da língua portuguesa e inglesa. O restante
deste relatório está organizado da seguinte forma: a Seção
\ref{sec:fundamentacao} apresenta a fundamentação teórica da cifra e
as decisões de implementação da Parte I; a Seção
\ref{sec:experimental} descreve o ambiente experimental, a
metodologia do ataque da Parte II e os resultados intermediários; a
Seção \ref{sec:conclusoes} apresenta as
conclusões técnicas; e a bibliografia lista as referências
utilizadas.

\section{Fundamentação Teórica}
\label{sec:fundamentacao}

\subsection{A cifra de Vigenère}

A cifra de Vigenère desloca cada letra do texto claro por um valor
que depende da posição correspondente de uma chave, repetida
ciclicamente sobre a mensagem. Para uma letra do texto claro $P_i$ e
a letra da chave na posição $i \bmod L$ ($L$ = tamanho da chave),
convertidas para valores numéricos de 0 a 25:

\begin{align}
\text{Cifração:}\quad   & C_i = (P_i + K_{i \bmod L}) \bmod 26 \\
\text{Decifração:}\quad & P_i = (C_i - K_{i \bmod L} + 26) \bmod 26
\end{align}

Diferentemente da cifra de César (deslocamento único), a chave do
Vigenère é uma sequência, gerando um espaço de chaves de tamanho
$26^L$ --- para $L = 14$, por exemplo, aproximadamente $2^{66}$
chaves possíveis, tornando a busca exaustiva inviável. Essa aparente
robustez não impede, contudo, o ataque estatístico descrito na Seção
\ref{sec:experimental}, precisamente porque a chave se repete
ciclicamente \cite{fundamentos}.

\subsection{Tratamento do alfabeto}

O grupo adotou a seguinte política, comum a toda a implementação:

\begin{itemize}
\item Alfabeto de cifração: A--Z (ASCII), sem distinção entre
  maiúsculas e minúsculas para o cálculo do deslocamento.
\item O caso original (maiúscula/minúscula) de cada letra é
  preservado no criptograma.
\item Espaços, números, pontuação e demais caracteres fora de A--Z
  são copiados sem alteração e \textbf{não consomem posição da
  chave}.
\item A chave é validada: deve ser não vazia e conter apenas letras
  A--Z.
\end{itemize}

\subsection{Fundamentos da criptoanálise (Índice de Coincidência)}

O ataque de recuperação de chave apoia-se no Índice de Coincidência
(IC), medida proposta por Friedman para quantificar a concentração
estatística das letras \cite{friedman}:

\begin{equation}
IC = \frac{\sum_i n_i (n_i - 1)}{N (N - 1)}
\end{equation}

onde $n_i$ é a frequência da letra $i$ no texto e $N$ o tamanho do
texto. Um texto em linguagem natural tem IC próximo de
$0{,}065$--$0{,}075$; um texto com distribuição uniforme de letras
(como um criptograma de Vigenère mal alinhado) tem IC próximo de
$1/26 \approx 0{,}038$. Essa diferença é a base para estimar o
tamanho da chave (detalhado na Seção \ref{sec:experimental}).

\section{Ambiente Experimental e Análise de Resultados}
\label{sec:experimental}

\subsection{Descrição do Cenário}

\begin{itemize}
\item \textbf{Linguagem e ferramentas:} C++17, biblioteca padrão,
  CMake 3.22.1 e testes executáveis próprios, sem bibliotecas de cifra
  ou criptoanálise pronta.
\item \textbf{Ambiente de compilação/execução:} Ubuntu 22.04, com
  compilador g++ 11.4.0. A validação desta versão foi feita localmente
  em modo de linha de comando.
\item \textbf{Dados experimentais:} dois criptogramas fornecidos pela
  disciplina --- um gerado a partir de texto em português, outro a
  partir de texto em inglês, cifrados com chaves diferentes.
  \pendente{adicionar os arquivos oficiais à pasta \texttt{textos/}}
\item \textbf{Estrutura do código:} cifrador
  (\texttt{parte1\_cifrador/}) implementado como biblioteca
  reutilizável, com interface de linha de comando, testes
  automatizados e benchmark de desempenho; módulo de criptoanálise
  (\texttt{parte2\_criptoanalise/}) com estimador de tamanho da chave
  e testes próprios.
\end{itemize}

\subsection{Análise de Resultados}

Metodologia do ataque de recuperação de chave, alinhada à
fundamentação da Seção \ref{sec:fundamentacao}:

\begin{enumerate}
\item \textbf{Estimativa do tamanho da chave}: para cada tamanho
  candidato $k$, o criptograma é dividido em $k$ subconjuntos
  (posições $0, k, 2k, \dots$; $1, k{+}1, 2k{+}1, \dots$; etc.) e o
  IC médio dos subconjuntos válidos é calculado. Os tamanhos são
  ordenados pelo maior IC médio; quando a diferença é inferior à
  tolerância de $0{,}003$, o menor tamanho é priorizado, reduzindo
  falsos candidatos que sejam múltiplos da chave real.
\item \textbf{Recuperação de cada letra da chave}: com o tamanho $k$
  estimado, cada subconjunto corresponde a um deslocamento fixo
  (equivalente a uma Cifra de César). Calcula-se a frequência
  relativa de cada letra do subconjunto e testam-se os 26
  deslocamentos possíveis, comparando a distribuição observada com a
  distribuição esperada do idioma \cite{frequencias}; o deslocamento de maior
  correlação será adotado como candidato à letra da chave.
\item \textbf{Reconstrução e validação}: as letras recuperadas
  formarão a chave candidata, que será usada para decifrar o criptograma
  completo. O resultado será avaliado quanto à legibilidade; se
  necessário, letras específicas são refinadas manualmente.
\end{enumerate}

Os componentes disponíveis foram validados separadamente. O cifrador
passou pelos nove grupos de testes existentes, incluindo o vetor
clássico, preservação de caixa, caracteres não alfabéticos, chave
inválida e 500 casos determinísticos. O estimador identificou como
primeira opção os tamanhos corretos 5, 7, 11 e 16 nos quatro
criptogramas de teste. Esses resultados verificam os módulos já
implementados, mas ainda não equivalem à validação do ataque completo.

\begin{table}[htbp]
\caption{Resultados obtidos \pendente{preencher após execução da Parte II}}
\label{tab:resultados}
\centering
\begin{tabular}{|l|l|l|l|}
\hline
\textbf{Criptograma} & \textbf{Idioma} & \textbf{Tam. chave} & \textbf{Legível?} \\
\hline
Criptograma 1 & Português & --- & --- \\
\hline
Criptograma 2 & Inglês & --- & --- \\
\hline
\end{tabular}
\end{table}

\textbf{Ajustes e problemas encontrados:} tamanhos múltiplos da chave
podem apresentar IC elevado; por isso, o estimador introduz uma
tolerância estatística e favorece o menor candidato em resultados
próximos. A recuperação das letras, a decifração e a avaliação de
legibilidade dependem da integração ainda pendente.

\section{Conclusões}
\label{sec:conclusoes}

Os resultados intermediários demonstram que um espaço de chaves grande
($26^L$) não é
garantia de segurança quando a chave se repete ciclicamente sobre o
texto --- a mesma vulnerabilidade que o Índice de Coincidência
explora. Os testes confirmam a correção da cifração e mostram que o IC
é capaz de priorizar o tamanho verdadeiro nas amostras disponíveis.
Ainda é necessário integrar a análise de frequências e avaliar os dois
criptogramas oficiais antes de concluir sobre a recuperação das chaves.
Como contraponto teórico, uma chave do mesmo tamanho do
texto, gerada aleatoriamente e usada uma única vez (\emph{One-Time
Pad}), tornaria esse ataque estatístico ineficaz --- a fragilidade
está na repetição da chave, não na estrutura da cifra em si.

\begin{thebibliography}{9}
\bibitem{roteiro}
Roteiro do Trabalho de Implementação 1 --- CIC0201, Profa. Priscila Solis, UnB, 2026.

\bibitem{fundamentos}
Slides de Criptografia Fundamentos (\textit{lec1-2026}, \textit{lec2-total}, \textit{introduction-part2}) --- CIC0201, UnB, 2026.

\bibitem{friedman}
W. F. Friedman, \textit{The Index of Coincidence and Its Applications
in Cryptography}. Riverbank Laboratories, 1922.

\bibitem{frequencias}
Wikipédia, ``Frequência de letras.'' [Online]. Disponível em:
\url{https://pt.wikipedia.org/wiki/Frequencia_de_letras}.
Acesso em: 31 ago. 2026.
\end{thebibliography}

\end{document}
